\chapter{Compile-Time Evaluation}
\label{chap:comptime}

Compile-time evaluation is the most distinctive feature of \epl{}. It allows a
programmer to write ordinary \epl{} expressions that the compiler executes while
compiling the program. The result is inserted back into the program as a
constant. This chapter explains why the feature exists, how purity is enforced,
how the evaluator represents values and memory, and how pure functions are
called during compilation.

\section{Purpose}

Without compile-time evaluation, a programmer who wants a precomputed table must
either write the table by hand, generate it with an external script, or compute
it at run time. Each option has a drawback. Hand-written tables are error-prone.
External generators add another tool and another language to the build process.
Run-time computation repeats work that could have been done once during
compilation.

\texttt{comptime} provides another option. The programmer writes a normal
expression. The compiler type checks it, proves that it is valid for
compile-time execution, interprets it, and replaces it with a constant. The
example in Listing \ref{lst:comptime-primes-full} computes the first twenty odd
prime numbers during compilation and then prints them at run time.

\begin{lstlisting}[language=EPL,caption={Compile-time prime table example},label={lst:comptime-primes-full}]
fn main() -> i32 {
	let primes = comptime {
		let lut: [i32; 20];
		let ready = 0;
		let next = 3;
		while ready < 20 {
			while !is_prime(next) {
				next += 2;
			}
			lut[ready as u64] = next;
			ready += 1;
			next += 2;
		}
		lut
	};
	for i in 0u64..20 {
		printf("%i ", primes[i]);
	}
	printf("\n");
	0
}
\end{lstlisting}

The important point is that the loop in the \texttt{comptime} block does not
remain in the generated run-time program. It is executed by the compiler, and
the result is a constant array in \irtree{}.

\section{Purity Boundary}

Executing arbitrary program code inside the compiler would be unsafe and
non-deterministic. A compile-time expression should not read run-time variables,
follow raw pointers, call external functions, or depend on the state of the
future process. \epl{} therefore defines a purity boundary.

A \texttt{comptime} expression may use:

\begin{itemize}
  \item numeric and boolean literals;
  \item arithmetic, comparison, boolean operations, and integer casts;
  \item blocks, conditionals, and loops;
  \item local variables declared inside the compile-time expression;
  \item assignment to those local variables;
  \item arrays and structs;
  \item \texttt{break} and \texttt{continue} for loops inside the expression;
  \item calls to functions annotated with \texttt{@pure}.
\end{itemize}

It may not use:

\begin{itemize}
  \item string literals;
  \item surrounding run-time variables or function arguments;
  \item pointer address-of or dereference operations;
  \item impure function calls;
  \item \texttt{return} from the enclosing function;
  \item loops or variables outside the compile-time expression's context.
\end{itemize}

This purity boundary is enforced after lowering to \irtree{}. The checker walks
the expression tree and validates each operation. Checking \irtree{} instead of
the AST is important because the tree has resolved IDs and types. The checker
can distinguish a variable declared inside the compile-time expression from a
variable in an enclosing run-time scope.

\section{Pure Functions}

The \texttt{@pure} annotation marks a function that is valid to call from
compile-time code. A pure function must have a body. It may use its own
arguments, local variables, local mutation, ordinary control flow, and calls to
other pure functions. It may not call impure functions, use strings, take
addresses, or dereference pointers.

Listing \ref{lst:pure-is-prime} shows a pure function from the examples.

\begin{lstlisting}[language=EPL,caption={Pure function used by comptime},label={lst:pure-is-prime}]
@pure
fn is_prime(x: i32) -> bool {
	if x <= 1 {
		return false;
	}
	if x == 2 {
		return true;
	}
	if x % 2 == 0 {
		return false;
	}
	let i = 3;
	while i <= x / 2 {
		if x % i == 0 {
			return false;
		}
		i += 2;
	}
	true
}
\end{lstlisting}

This function uses local mutation and early returns, but it does not depend on
run-time state. Therefore the compile-time evaluator can call it with constant
arguments.

\section{Evaluator Structure}

The evaluator is implemented in \texttt{src/ir\_tree/evaluator.rs}. It evaluates
\irtree{} expressions and returns an \texttt{ir\_tree::Constant}. The public
entry point for compile-time expressions is \texttt{eval\_comptime\_expr}. Pure
function calls are handled by a separate internal function,
\texttt{eval\_pure\_function}.

The evaluator uses a recursive expression interpreter. It pattern matches on
\texttt{ExprKind}. Constants evaluate to themselves. Loads read from compile-time
memory. Stores write to compile-time memory and return unit. Arithmetic and
comparison call helper routines. Conditionals evaluate the condition and then
one branch. Loops repeatedly evaluate the body until a matching break is
encountered. Function calls evaluate their arguments and then interpret the body
of the pure callee.

\section{Control-Flow Signals}

The evaluator needs to support \texttt{return}, \texttt{break}, and
\texttt{continue}. These are not ordinary values because they change control
flow. The implementation represents them as internal evaluation errors:
\texttt{Return(Constant)}, \texttt{Break(LoopId, Constant)}, and
\texttt{Continue(LoopId)}.

This design is similar to using exceptions internally in an interpreter. When a
loop body produces a \texttt{Break} with the loop's own ID, the loop consumes the
signal and produces the break value. When it produces a \texttt{Continue} with
the loop's own ID, the loop starts the next iteration. If the signal belongs to
an outer loop or a return, it propagates outward.

At the top level of a \texttt{comptime} expression, \texttt{return},
out-of-context \texttt{break}, and out-of-context \texttt{continue} should be
impossible because the purity checker rejects them. If one appears there, it is
a compiler bug.

\section{Constant Memory}

Local mutation inside \texttt{comptime} requires memory. The evaluator cannot
only store a map from variables to high-level constants because field and array
element assignment require partial updates. For example:

\begin{lstlisting}[language=EPL,caption={Partial compile-time mutation}]
let lut: [i32; 20];
lut[ready as u64] = next;
\end{lstlisting}

The evaluator represents each local variable as a byte buffer sized according to
the same type layout used by the compiler. A compile-time place is represented
as a variable ID plus a byte offset. Loading converts bytes back into a constant.
Storing converts a constant into bytes and writes them into the buffer.

The helper module \texttt{const\_abi.rs} performs conversion between constants
and byte buffers. It handles integers, booleans, arrays, structs, unit, and
undefined values. For structs, it respects field offsets and padding. For arrays,
it serializes elements in order using the element layout.

This design keeps compile-time mutation consistent with the run-time layout used
later by lower IR and LLVM generation.
The byte-addressed model was also chosen so that pointer support could be added
to compile-time evaluation later without replacing the evaluator's storage
abstraction. Such support would be restricted to pointers into local
compile-time variables, with checked pointer arithmetic, explicit pointer casts,
and dereferences validated against the known local memory ranges.

\section{Arithmetic and Casts}

The evaluator's arithmetic helper supports all integer types. It uses wrapping
operations for addition, subtraction, multiplication, division, and remainder.
Comparison supports booleans and integers. Cast evaluation supports
integer-to-integer casts by using Rust casts to model truncation and extension.

Pointer constants are not supported by compile-time evaluation. This is an
intentional restriction. A pointer value would either point into compile-time
memory, which has no valid run-time address, or into run-time memory, which does
not exist yet. Rejecting pointer operations keeps the compile-time value model
simple and deterministic.

\section{Case Study: Struct Computation}

The struct example demonstrates that compile-time evaluation is not limited to
primitive integers. The program defines a \texttt{Vec2} struct and a pure
\texttt{vec2\_add} function. At run time it prints ordinary struct values. At
compile time it also evaluates:

\begin{lstlisting}[language=EPL,caption={Compile-time struct expression}]
vec2_print(comptime vec2_add(.{ x: 1, y: 2 }, .{ y: 10, x: 20 }));
\end{lstlisting}

The \irtree{} snapshot shows that this call becomes a constant struct:

\begin{lstlisting}[caption={Struct constant after compile-time evaluation}]
FUNCTION_CALL("vec2_print") TYPE=unit
|   CONST {21, 12} TYPE=StructId(0)
\end{lstlisting}

The field order in the initializer does not matter. Semantic analysis maps
fields by name, and constant construction stores them in declaration order.

\section{Benefits and Tradeoffs}

The benefit of this design is that compile-time evaluation reuses the normal
language. There is no separate macro language, no string-based code generation,
and no second parser. The compiler can type check code once, then either lower
it to run time or interpret it during compilation.

The tradeoff is that purity rules must be conservative. Some operations that
could theoretically be made safe are rejected. For example, the evaluator could
eventually support references into compile-time-only memory, but the current
pointer model does not distinguish compile-time pointers from run-time pointers.
Rejecting pointer operations is simpler and safer for this implementation.

The current evaluator is also not optimized for speed. It repeatedly traverses
the module to find and replace compile-time expressions. The documentation notes
that this is inefficient. For the project examples it is sufficient, but a
larger compiler would use a more direct worklist or integrate evaluation into
semantic lowering.
